% =============================================================================
% Module 11: Solutions to Exercises
% =============================================================================
% This file is conditionally included based on \ifsolutions
% =============================================================================

\section{Solutions to Exercises}
\label{sec:statistics_solutions}

\noindent
Full worked solutions are also available in the Mathematica script
\solnlink{11_Lensing_Statistics/problems\_11.wl}{problems\_11.wl},
which verifies each answer symbolically.

% ----- Exercise 11.1 -----
\subsection*{Exercise~\ref{ex:point_cross}: Point-Mass Magnification Cross-Section}

\textbf{(a)} The total magnification of a point-mass lens is
$\mu_p(y) = (y^2+2)/\big(y\sqrt{y^2+4}\big)$, a monotonically decreasing
function of $y$.  Solving $\mu_p(y) = \mu$ for $y^2$: square both sides,
\[
    \mu^2 = \frac{(y^2+2)^2}{y^2(y^2+4)}
    \;\Longrightarrow\;
    \mu^2\,y^2(y^2+4) = (y^2+2)^2 .
\]
Writing $u = y^2$ and expanding gives
$\mu^2 u^2 + 4\mu^2 u = u^2 + 4u + 4$, i.e.
$(\mu^2-1)u^2 + 4(\mu^2-1)u - 4 = 0$.  The positive root is
\[
    u = y^2 = -2 + \frac{2\mu}{\sqrt{\mu^2-1}}
    = 2\left(\frac{\mu}{\sqrt{\mu^2-1}} - 1\right),
\]
which is Eq.~\eqref{eq:point_ymu}.  Since $\mu_p$ decreases with $y$, the
condition $\mu_p > \mu$ is $y < y(\mu)$, so the cross-section is
$\sigma(\mu) = \pi\thetaE^2\,y^2(\mu)$. \checkmark

\textbf{(b)} For $\mu \gg 1$, expand
$\mu/\sqrt{\mu^2-1} = (1-\mu^{-2})^{-1/2} = 1 + \tfrac12\mu^{-2} + \cdots$,
so $y^2(\mu) = 2(\mu/\sqrt{\mu^2-1} - 1) \to \mu^{-2}$.  Equivalently
$\lim_{\mu\to\infty}\mu^2\,y^2(\mu) = 1$, hence $y^2 \sim 1/\mu^2$ and
$\sigma \propto \mu^{-2}$. \checkmark  Since the differential probability
is $p(\mu)\,d\mu \propto -\,(d\sigma/d\mu)\,d\mu$, this cross-section is
the integral form of the universal $p(\mu)\propto\mu^{-3}$ tail.

% ----- Exercise 11.2 -----
\subsection*{Exercise~\ref{ex:sis_stats}: SIS Flux Ratio and Magnification}

For an SIS with $y = \beta/\thetaE < 1$ the two images have signed
magnifications $\mu_\pm = 1 \pm \thetaE/\beta = 1 \pm 1/y$.  The trailing
image ($\mu_-$) has $|\mu_-| = 1/y - 1$ (it is a saddle, negative parity).
The total magnification is
\[
    \mu = |\mu_+| + |\mu_-| = \left(1+\tfrac1y\right)
    + \left(\tfrac1y-1\right) = \frac{2}{y},
\]
and the flux ratio is
\[
    r = \frac{|\mu_+|}{|\mu_-|}
    = \frac{1+1/y}{1/y-1} = \frac{1+y}{1-y}. \quad\checkmark
\]
Since $r$ increases from $1$ (at $y=0$) to $\infty$ (at $y=1$), requiring a
flux ratio smaller than $r$ restricts the source to $y < (r-1)/(r+1) < 1$,
so the cross-section is bounded by the full multiple-imaging disk,
$\sigma \le \pi\thetaE^2$. \checkmark

% ----- Exercise 11.3 -----
\subsection*{Exercise~\ref{ex:bias_null}: The $\beta=1$ Null of Magnification Bias}

With $N_0(>S) = A\,S^{-\beta}$, Eq.~\eqref{eq:mag_bias} gives
\[
    N(>S) = \frac{1}{\langle\mu\rangle}\int d\mu\; p(\mu)\,
    A\left(\frac{S}{\mu}\right)^{-\beta}
    = \frac{A\,S^{-\beta}}{\langle\mu\rangle}
      \int d\mu\; p(\mu)\,\mu^{\beta}
    = N_0(>S)\,\frac{\langle\mu^{\beta}\rangle}{\langle\mu\rangle}.
\]
At $\beta = 1$ the numerator is
$\int \mu\,p(\mu)\,d\mu = \langle\mu\rangle$ by the definition of the mean
magnification, so $N(>S) = N_0(>S)$ for \emph{any} $p(\mu)$. \checkmark

\textbf{Physical reason:} magnification does two competing things.  It
brightens each source by $\mu$ (raising it above threshold), but it also
stretches the sky by $\mu$, diluting the surface density of sources by
$1/\mu$.  For counts with logarithmic slope $\beta = 1$ (equal flux per
logarithmic interval) these effects cancel exactly; steeper counts
($\beta>1$) are net-boosted, flatter counts ($\beta<1$) net-depleted.
